\chapter{Cơ sở lý thuyết}




\section{Tổng quan về mã độc PE trên hệ điều hành Windows}




\subsection{Tổng quan về mã độc (malware)}

Mã độc (malware - malicious software) là thuật ngữ dùng để
chỉ các chương trình phần mềm được thiết kế với mục đích
gây hại cho hệ thống máy tính, đánh cắp dữ liệu, gián điệp
người dùng, hoặc thực hiện các hành vi trái phép khác mà
không có sự cho phép của chủ sở hữu hệ thống. Mã độc có
thể lây lan qua nhiều con đường khác nhau như email,
phần mềm giả mạo, thiết bị lưu trữ di động hoặc qua
mạng Internet.

Các nhà nghiên cứu chia mã độc thành nhiều loại khác nhau,
và mỗi loại mã độc lại có thể được phân loại chi tiết hơn
(chia thành các phân nhóm nhỏ hơn). Một số loại mã độc phổ
biến trong hệ thống máy tính và hệ thống mạng bao gồm: \cite{wikimalware}

\begin{enumerate}
    \item \textbf{Virus máy tính (Computer virus):}
	Lây nhiễm bằng cách chèn mã độc vào các tập tin hợp lệ
	(vật chủ) và phát tán khi người dùng mở các tập tin đó.

	\item \textbf{Sâu (Worm):} Loại mã độc có khả năng
	tự nhân bản và lây lan nhanh qua mạng mà không cần
	đến tập tin chủ (vật chủ) như virus.

    \item \textbf{Trojan:} Loại mã độc ngụy trang dưới
	dạng phần mềm hợp lệ để đánh lừa người dùng cài đặt.

    \item \textbf{Mã độc tống tiền (Ransomware):}
	Mã hóa dữ liệu của nạn nhân và đòi tiền chuộc để khôi phục.

    \item \textbf{Phần mềm gián điệp (Spyware):}
	Theo dõi và ghi lại hành vi người dùng để đánh cắp thông
	tin nhạy cảm. Một trong những dạng spyware đầu tiên là
	keylogger - ghi lại các thao tác gõ phím của người dùng.
	Một số phần mềm gián điệp hiện đại còn khai thác thông
	tin kênh bên (side-channel) từ các đặc tính vật lý của
	hệ thống (như điện năng tiêu thụ, nhiệt độ, thời gian
	thực thi\dots) nhằm phục vụ cho việc phân tích, giải mã
	thông tin mã hóa, bẻ khóa hệ thống, truy cập trái
	phép\dots

    \item \textbf{Rootkit:} Loại mã độc ẩn mình sâu trong
	hệ thống, che giấu các tiến trình hoặc tệp tin liên
	quan để tránh bị phát hiện. Rootkit có thể làm được
	điều này là nhờ nó can thiệp sâu vào kernel hoặc
	các thành phần hệ thống quan trọng, cho phép nó
	kiểm soát hệ thống cũng như các phần mềm bảo mật.
\end{enumerate}

Mã độc có thể được chứa trong nhiều loại tệp tin, thậm chí cả PDF.
Trên hệ điều hành Windows, một trong những loại tệp tin chứa mã độc
phổ biến hơn cả là định dạng PE (Portable Executable) - định dạng
tập tin thực thi tiêu chuẩn của Windows, thường có phần mở rộng
.exe, .dll, .sys... Mã độc PE lợi dụng cấu trúc đặc trưng của
định dạng này, cũng như những lỗ hổng trong cách thức hệ điều hành
Windows nạp (load) chúng vào bộ nhớ, từ đó chèn mã độc vào các
vùng nhớ, giấu mã hoặc thực hiện các hành vi tinh vi như tải
xuống mã độc bổ sung, thực thi theo điều kiện, hoặc chống phân
tích ngược\dots

Mã độc hiện đại thường là một phần trong các chiến dịch tấn công
có chủ đích (APT - Advanced Persistent Threat). Các chiến dịch
này được tổ chức bài bản, kéo dài trong thời gian dài và thường
nhắm vào các tổ chức, doanh nghiệp hoặc cơ quan chính phủ để
đánh cắp thông tin, gây rối hoạt động hoặc phá hoại. \cite{wikiapt}

Hiện nay, thế giới đang chứng kiến sự gia tăng mạnh mẽ cả về số
lượng lẫn mức độ tinh vi của các loại mã độc. Theo báo cáo của
AV-TEST (2024), \cite{avtestmalwarestats} trung bình mỗi ngày có hơn 450.000 mẫu mã độc
mới được phát hiện trên toàn cầu, cho thấy mức độ phát triển
nhanh chóng và khó kiểm soát của các mối đe dọa này. Đặc biệt,
mã độc dạng PE (Portable Executable) vẫn chiếm tỷ lệ lớn trong
các cuộc tấn công nhắm vào hệ điều hành Windows, nền tảng phổ
biến trong cả môi trường cá nhân lẫn doanh nghiệp. Do khả năng
thực thi dễ dàng, linh hoạt và tích hợp sâu vào hệ thống, mã
độc PE có thể lợi dụng nhiều kỹ thuật để tránh bị phát hiện,
đồng thời phát tán nhanh chóng trong mạng nội bộ của tổ chức.
Điều này đặt ra rủi ro nghiêm trọng về mất mát dữ liệu, gián
đoạn hoạt động kinh doanh và thậm chí là ảnh hưởng đến an ninh
quốc gia nếu các cơ quan chính phủ bị xâm nhập.

Trong bối cảnh chuyển đổi số đang diễn ra mạnh mẽ tại Việt Nam
cũng như nhiều quốc gia trên thế giới theo định hướng Cách mạng
công nghiệp 4.0, các hệ thống thông tin ngày càng trở nên phức
tạp, kết nối chặt chẽ và phụ thuộc nhiều hơn vào công nghệ.
Điều này đồng nghĩa với việc mở rộng bề mặt tấn công (attack
surface) cho các đối tượng xấu và làm gia tăng nguy cơ bị mã độc
tấn công, đặc biệt trong các tổ chức trọng yếu như ngân hàng,
y tế, cơ quan nhà nước hay doanh nghiệp hạ tầng. Chính vì vậy,
nhu cầu xây dựng và triển khai các phương pháp phát hiện và
ngăn chặn mã độc hiệu quả, trong đó có mã độc PE, đang trở
thành một vấn đề cấp bách, không chỉ đối với các tổ chức
chuyên về an ninh mạng mà còn đối với toàn bộ hệ sinh thái
số quốc gia.













\subsection{Tổng quan về các phương pháp phát hiện và ngăn chặn mã độc}

Mã độc ngày nay không chỉ đơn giản là những đoạn mã phá hoại mà
đã phát triển thành các công cụ tấn công tinh vi, được sử dụng
trong các chiến dịch APT, ransomware, gián điệp công nghiệp và
thậm chí là tấn công vào cơ sở hạ tầng quốc gia. Trước thực trạng
đó, không một giải pháp đơn lẻ nào có thể đảm bảo an toàn tuyệt
đối, mà cần đến một chiến lược an ninh mạng nhiều lớp
(multi-layered security), kết hợp các phương pháp phát hiện và
ngăn chặn từ nhiều góc độ, chẳng hạn như:

\begin{itemize}
	\item \textbf{Lớp bảo vệ mạng:} Phát hiện mã độc qua lưu
	lượng mạng, sandbox, IDS/IPS.

	\item \textbf{Lớp cổng kết nối:} Ngăn chặn mã độc qua email,
	web, tệp tải về.

	\item \textbf{Lớp endpoint:} Phát hiện và phản ứng trực tiếp
	trên thiết bị đầu cuối - nơi mã độc thực sự được thực thi.

	\item \textbf{Lớp ứng dụng/cloud:} Bảo vệ dịch vụ, dữ liệu
	và nền tảng trong môi trường đám mây.

	\item \textbf{Lớp phân tích và phản ứng:} Bao gồm Threat
	Intelligence, SOC, SIEM, Threat Hunting\dots
\end{itemize}

Trong đó, việc bảo vệ endpoint (thiết bị đầu cuối) \cite{wikiendpointsecurity} đóng vai
trò đặc biệt quan trọng bởi đây là nơi mã độc thực sự tác động
đến người dùng, dữ liệu và tài sản số của tổ chức. Các thiết bị
đầu cuối như máy tính cá nhân, máy chủ, laptop nhân viên, thiết
bị IoT\dots là đích đến cuối cùng của phần lớn mã độc. Dù tấn
công bắt đầu từ email, trình duyệt hay thiết bị USB, mã độc vẫn
cần được thực thi trên endpoint để gây hại. Do đó, bảo vệ
endpoint là tuyến phòng thủ cuối cùng và quan trọng nhất, giúp
tổ chức phát hiện sớm các hành vi đáng ngờ ngay tại thiết bị,
ngăn chặn mã độc trước khi nó thực hiện hành vi phá hoại hoặc
lây lan, cũng như có khả năng phản ứng kịp thời, cách ly thiết
bị, thu thập bằng chứng phục vụ điều tra. Một giải pháp phát
hiện mã độc ngay tại máy endpoint của người dùng có thể được
coi là rất hữu ích và hiệu quả.

Trên cơ sở đó, việc khóa luận tập trung phát triển tính năng
agent phát hiện mã độc chạy trên máy endpoint là hoàn toàn phù hợp
với nhu cầu thực tiễn. Thay vì chỉ dừng lại ở việc xây dựng mô hình
học máy thuần túy, khóa luận hướng đến tích hợp mô hình vào một hệ
thống có thể vận hành thực sự trên môi trường Windows - nơi mã độc
PE thực sự hoạt động và gây hại. Tuy nhiên, việc xây dựng một giải pháp bảo vệ endpoint hoàn chỉnh
đòi hỏi sự kết hợp của nhiều lớp kỹ thuật khác nhau, từ phân tích
và phân loại file cho đến giám sát hành vi thời gian thực ở tầng hệ
điều hành. Để giải quyết hiệu quả bài toán phân loại, cần có một
pipeline học máy đủ mạnh, được xây dựng trên bộ đặc trưng phong phú
và có ý nghĩa ngữ nghĩa. Bên cạnh đặc trưng cấu trúc thông thường
của file PE, việc bổ sung thông tin về \textit{năng lực} (capability)
mà file có thể thực hiện - chẳng hạn khả năng mã hóa dữ liệu, kết
nối mạng, leo thang đặc quyền hay can thiệp vào tiến trình hệ thống
- mang lại góc nhìn ngữ nghĩa sâu hơn, giúp mô hình phân biệt mã
độc hiệu quả hơn so với chỉ dựa vào đặc trưng thống kê đơn thuần.
Chi tiết việc xây dựng và tích hơp pipeline sẽ được trình bày,
làm rõ trong Chương 2 và Chương 3.






\subsection{Tổng quan về x86/amd64 Assembly}

Trong lĩnh vực an toàn thông tin, đặc biệt là trong phân tích và
phát hiện mã độc trên hệ điều hành Windows, kiến thức nền tảng về
kiến trúc vi xử lý x86/amd64 và định dạng tệp thực thi PE
(Portable Executable) đóng vai trò rất quan trọng. Đây là hai
yếu tố kỹ thuật cốt lõi giúp hiểu rõ cách mã độc hoạt động,
cách chúng lẩn tránh các cơ chế bảo vệ, cũng như cách các
công cụ phân tích có thể phát hiện và vô hiệu hóa chúng.

Windows hiện vẫn là hệ điều hành phổ biến nhất cho người dùng
cá nhân và doanh nghiệp. Các mã độc nhắm vào nền tảng này - đặc
biệt là dạng mã độc PE - tiếp tục chiếm tỷ lệ lớn trong các
chiến dịch tấn công thực tế. Với khả năng can thiệp trực tiếp
vào bộ nhớ, ẩn mình trong tiến trình hợp pháp, hoặc khai thác
các kỹ thuật cấp thấp, mã độc PE đòi hỏi người phòng thủ phải
nắm vững kiến thức kỹ thuật ở tầng thấp (low-level), cụ thể là:

\begin{itemize}
    \item Cấu trúc và cách thức hoạt động của mã máy (Assembly)
	trên nền kiến trúc x86/amd64. Hiện nay, đây vẫn là kiến trúc
	tập lệnh (ISA) phổ biến nhất được sử dụng bởi các thiết bị
	Windows cũng như các tệp PE.
    \item Định dạng và đặc điểm nội tại của các tệp PE.
\end{itemize}

Việc hiểu được cách các mã máy được sinh ra, tổ chức trong
một tệp PE, và cuối cùng được nạp vào bộ nhớ và thực thi,
là nền tảng không thể thiếu cho các chuyên gia làm việc trong
các lĩnh vực như phân tích mã độc (malware analysis), kỹ thuật
dịch ngược (reverse engineering) và xây dựng hệ thống phát
hiện mã độc bằng học máy (ML-based malware detection).
Khóa luận này sẽ trình bày chi tiết về hai chủ đề chính:
kiến trúc x86/amd64 và định dạng tệp PE, trong đó đặc biệt
nhấn mạnh tới các điểm dễ bị lợi dụng bởi mã độc cũng như
các đặc trưng hữu ích cho việc phát hiện và phân tích chúng.

Hai kiến trúc x86 (32-bit) và amd64 (64-bit, còn gọi là x86-64)
là nền tảng xử lý phổ biến nhất trên máy tính cá nhân và máy chủ
sử dụng hệ điều hành Windows. Cả hai kiến trúc này đều thuộc họ
CISC (Complex Instruction Set Computing), với tập lệnh đa dạng
và phong phú, cho phép lập trình viên thao tác trực tiếp với
bộ nhớ, thanh ghi, và luồng điều khiển chương trình. \cite{intel-sdm} \cite{amd64-manual} \cite{ms-x64-calling}

\textbf{Về thanh ghi và cấu trúc bộ xử lý}, kiến trúc x86
cung cấp 8 thanh ghi 32-bit tổng quát: EAX, EBX, ECX, EDX,
ESI, EDI, EBP, ESP. Kiến trúc amd64 mở rộng số lượng
thanh ghi tổng quát lên 16, với kích thước 64-bit: RAX,
RBX, RCX, RDX, RSI, RDI, RBP, RSP và 8 thanh ghi còn lại
R8-R15. Các thanh ghi ESP, RSP (Stack Pointer) và EIP,
RIP (Instruction Pointer) đóng vai trò quan trọng trong
điều khiển luồng thực thi, và cũng bị lợi dụng bởi một số
kỹ thuật tấn công phổ biến, chẳng hạn return-oriented
programming (ROP).

\textbf{Về chế độ vận hành}, hai kiến trúc đều cung cấp
các chế độ real mode, protected mode và long mode.
Trong real mode (chế độ thực), CPU không thực hiện bảo
vệ bộ nhớ cấp phần cứng (hardware-level memory protection) - tức
không gian địa chỉ luôn ánh xạ 1:1 với địa chỉ thực của bộ nhớ.
Trong khi đó, protected mode (chế độ bảo vệ) được sử dụng rộng
rãi trong x86, hỗ trợ phân vùng bộ nhớ và quyền truy cập,
không cho phép các chương trình đọc, ghi trực tiếp vào địa
chỉ thực của bộ nhớ mà phải sử dụng địa chỉ ảo (virtual addresses).
Long mode là chế độ chỉ có ở amd64, cho phép truy cập không gian
bộ nhớ lên tới 64 bit, sử dụng paging để bảo vệ vùng nhớ.

\textbf{Về tập lệnh và cú pháp}, hai kiến trúc đều có
tập lệnh hợp ngữ (assembly) rất đa dạng và có nhiều lệnh
(mnemonics) tương đồng, bao gồm thao tác số học (ADD,
SUB, INC, DEC), logic (AND, OR, XOR), đọc, ghi dữ liệu
(MOV, PUSH, POP), điều khiển luồng (CALL, JMP, RET, LOOP,
JZ, JNZ\dots). Cú pháp phổ biến của hợp ngữ trên Windows
là Intel syntax. Ngoài ra còn có AT\&T syntax. Cần lưu ý
rằng một đoạn mã máy có thể được dịch trở lại thành mã
hợp ngữ khác nhau do sự khác biệt về syntax hay các biến
thể của các trình hợp dịch (assembler) và dịch ngược
(disassembler).

Ngoài ra, trong phân tích mã độc và nghiên cứu file thực thi nhị phân (binary),
việc hiểu về quy ước hay giao ước gọi hàm (calling convention) cũng rất
quan trọng. Một calling convention quyết định cách các tham số được
truyền vào hàm, cách hàm đó lưu giá trị trả về, và cách
vùng nhớ stack được quản lý khi gọi hàm đó. Các calling
convention phổ biến được cho trong \ref{tab:calling-conventions}. Biết được
calling convention, người ta có thể xác định điểm vào
hàm, tham số, và luồng điều khiển khi xem disassembly - điều
cực kỳ quan trọng khi phân tích thủ công các tệp thực thi PE.
Hơn nữa, mã độc thường thao túng stack hoặc gọi các API hệ
thống một cách gián tiếp (indirect call, ví dụ dùng
\texttt{LoadLibrary} và \texttt{GetProcAddress} trên Windows),
khiến việc hiểu rõ calling convention là cần thiết để truy vết
hành vi thực sự của mã độc. Một số mã độc còn cố tình vi
phạm calling convention chuẩn để gây lỗi, gây khó hiểu (confuse)
cho các công cụ phân tích tĩnh như disassembler, công cụ dịch
ngược sang ngôn ngữ C (chẳng hạn Ghidra, IDA)\dots

Có thể thấy, các lệnh hợp ngữ chỉ thực hiện các thao tác đơn
giản, hơn nữa các calling convention chỉ là giao ước giữa
các hàm với nhau, hoàn toàn không bị bắt buộc hay bị kiểm
tra bởi CPU. Do đó, kiến trúc x86/amd64 tạo điều kiện cho
các cuộc tấn công khai thác lỗ hổng ở tầng thấp, chẳng hạn
buffer overflow, return-oriented programming (ROP) hay
shellcode injection. Nhiều mã độc được viết
bằng assembly tối giản, nhỏ gọn, khó phát hiện khi chưa
giải mã hoặc unpack. Việc hiểu và áp dụng hợp ngữ trên
Windows, đặc biệt là x86 và amd64, đóng vai trò quan
trọng trong phân tích mã độc, cũng như nghiên cứu giải
pháp phát hiện và ngăn chặn mã độc mới.

\begin{table}[h]
  \centering
  \begin{tabularx}{\textwidth}{|X|X|X|X|X|}
  \hline
  \textbf{Tên} & \textbf{Kiến trúc} & \textbf{Truyền tham số \textit{(kiểu số nguyên hoặc pointer)}} & \textbf{Trách nhiệm clean up stack} & \textbf{Ghi chú} \\
  \hline
  \texttt{cdecl} & x86 & Đẩy tham số lên stack (tham số cuối được đẩy trước) & Caller & Dùng phổ biến trong C \\
  \hline
  \texttt{stdcall} & x86 & Như \texttt{cdecl} & \textit{Callee} & Dùng trong WINAPI \\
  \hline
  \texttt{fastcall} & x86 & ECX, EDX, và đẩy các tham số còn lại lên stack như \texttt{cdecl} & \textit{Callee} & Nhanh hơn, thường thấy trong tối ưu (compiler optimization), song không phổ biến \\
  \hline
  \texttt{Microsoft x64} & amd64 & RCX, RDX, R8, R9, và đẩy các tham số còn lại lên stack như \texttt{cdecl} & Caller & Mặc định trên Windows 64-bit \\
  \hline
  \end{tabularx}
  \caption{Các calling convention phổ biến trên Windows}
  \label{tab:calling-conventions}
\end{table}








\subsection{Tổng quan về định dạng PE (Portable Executable) trên Windows}

Định dạng PE (Portable Executable) \cite{ms-pe-format} \cite{rickpefilestructure}
là định dạng chuẩn
cho các tệp thực thi trên hệ điều hành Windows, bao gồm \texttt{.exe},
\texttt{.dll}, \texttt{.sys}, và một số định dạng khác. PE là phần mở
rộng của định dạng COFF (Common Object File Format), được Microsoft
phát triển để hỗ trợ nạp mã máy vào bộ nhớ và thực thi một cách linh
hoạt. Mỗi tệp PE là một tổ chức chặt chẽ giữa mã máy (executable
code), dữ liệu (data), siêu dữ liệu (metadata), và thông tin cần thiết
để hệ điều hành có thể thực thi chương trình đúng ngữ cảnh.

Cấu trúc tổng quát của tệp PE gồm 5 bộ phận chính kế tiếp nhau:

\begin{itemize}
    \item \textbf{DOS Header} (còn gọi là MS-DOS Header): vùng có kích
        thước 64 bytes nhằm mục đích tương thích ngược với hệ điều hành
        MS-DOS. Nếu chạy tệp PE trong DOS, header này khiến DOS nhận
        diện tệp là một tệp thực thi bình thường và sẽ thực thi phần mã
        máy trong DOS stub. Đối với Windows NT, bộ phận này chỉ còn tồn
        tại vì lý do lịch sử. Đặc trưng của DOS Header là mở đầu bằng
        2 byte \texttt{"MZ"} (\texttt{0x5A4D}) - được gọi là
        \textit{magic number} - là dấu hiệu nhận diện tệp thực thi
        của DOS.

    \item \textbf{DOS Stub:} vùng chứa mã máy chạy được trong DOS.
        Thông thường vùng này in ra dòng thông báo \textit{``This
        program cannot be run in DOS mode''} rồi thoát chương trình.
        Phần này chỉ tồn tại vì lý do tương thích ngược; hệ điều hành
        Windows không thực thi mã máy trong phần này.

    \item \textbf{PE Header} (NT Header): vùng chứa các metadata -
        thông tin tổng thể về tập tin, chẳng hạn kiến trúc CPU của mã
        thực thi, số sections, và characteristics (đặc tính tập tin:
        executable, system file\dots). Đặc trưng của vùng này là mở đầu
        bằng một signature chứa 4 bytes \texttt{0x50450000}, tức
        \texttt{"PE{\textbackslash}0{\textbackslash}0"} theo ASCII.

    \item \textbf{Section Headers:} chứa thông tin từng section, bao
        gồm địa chỉ bắt đầu, kích thước và các thuộc tính của section.

    \item \textbf{Sections:} vùng chứa nội dung thực sự của các
        section, với vị trí và thuộc tính đã được định nghĩa trong
        Section Headers. Các sections điển hình bao gồm:
        \begin{itemize}
            \item \texttt{.text}: chứa mã thực thi.
            \item \texttt{.data}: biến toàn cục/biến tĩnh có thể đọc
                và ghi.
            \item \texttt{.rdata}: dữ liệu chỉ đọc.
            \item \texttt{.reloc}: bảng ánh xạ lại địa chỉ
                (relocations).
            \item \texttt{.rsrc}: tài nguyên (resources) như ảnh,
                icon, form\dots
        \end{itemize}
\end{itemize}

Kiến trúc x86/amd64 và định dạng PE trên Windows, tuy linh hoạt và
hiệu quả cho thực thi mã, nhưng cũng ẩn chứa nhiều điểm yếu có thể
bị mã độc lợi dụng để xâm nhập, ẩn mình và hoạt động âm thầm trong
hệ thống. Dưới đây liệt kê một số điểm yếu tiêu biểu:

\begin{itemize}
    \item \textbf{Bộ nhớ phẳng (flat) và chia sẻ (shared):} Trên
        Windows, tiến trình có thể truy cập toàn bộ không gian địa chỉ
        ảo của mình. Nếu không có cơ chế bảo vệ như DEP (Data Execution
        Prevention) hay ASLR (Address Space Layout Randomization), mã
        độc có thể thay đổi dữ liệu, ghi đè mã, hoặc thực thi từ vùng
        dữ liệu một cách dễ dàng.

    \item \textbf{Thiếu kiểm soát chặt chẽ giữa code và data:} Kết
        hợp với kỹ thuật code injection (tiêm mã vào tiến trình khác),
        mã độc có thể viết và thực thi mã ở bất kỳ đâu nếu có quyền
        phù hợp. Chẳng hạn, mã độc có thể che giấu payload bằng cách
        mã hóa, sau đó giải mã ra một vùng nhớ và đặt quyền thực thi
        trên vùng đó để kích hoạt mã độc.

    \item \textbf{Indirect calls/jumps:} Assembly cho phép gọi hàm
        hoặc nhảy đến địa chỉ được lưu trong thanh ghi hoặc bộ nhớ
        (ví dụ: \texttt{CALL [EAX]}, \texttt{JMP RAX}). Mã độc dùng
        kỹ thuật này để ẩn đích thực thi, gây khó khăn cho phân tích
        tĩnh và các công cụ dịch ngược.

    \item \textbf{Gọi API động:} Mã độc thường sử dụng
        \texttt{LoadLibrary}, \texttt{GetProcAddress}, resolve hashing,
        hoặc kỹ thuật gọi hàm theo địa chỉ offset để lẩn tránh các bộ
        lọc theo chuỗi API tĩnh.
\end{itemize}

Lợi dụng linh hoạt các điểm yếu trên, người ta đã phát triển rất
nhiều kỹ thuật khai thác nhắm vào tập tin PE và môi trường thực thi
Windows, chẳng hạn:

\begin{itemize}
    \item \textbf{Section spoofing:} Tạo các section giả với tên hợp
        lệ như \texttt{.text} hoặc \texttt{.rdata} nhưng chứa mã độc,
        gây nhầm lẫn cho công cụ phân tích; hoặc tạo section tùy ý
        chứa mã độc (\texttt{.xyz}, \texttt{.UPX}, \texttt{.bss},\dots).

    \item \textbf{PE header corruption:} Thay đổi hoặc làm hỏng PE
        Header nhằm làm gián đoạn công cụ phân tích hoặc antivirus.

    \item \textbf{Entry point manipulation:} Thay đổi Entry Point để
        trỏ vào mã độc thay vì hàm \texttt{main} hoặc \texttt{WinMain}
		thông thường.

    \item \textbf{Import table spoofing:} Tạo bảng IAT (Import Address
        Table) giả, hoặc mã hóa các chuỗi API để vượt qua signature
        detection.

    \item \textbf{Packing/obfuscation:} Sử dụng packer, crypter để mã
        hóa toàn bộ phần \texttt{.text}, chỉ giải mã khi chạy, từ đó
        vô hiệu hóa phân tích tĩnh thông thường. Một số kỹ thuật làm
        nhiễu trong Relocation và Exception Table cũng thường được sử
        dụng.

    \item \textbf{Process hollowing/doppelgänging:} Tạo một tiến trình
        hợp pháp (như \texttt{notepad.exe}), rồi ghi đè bộ nhớ của
        tiến trình đó bằng mã độc, giúp chạy dưới tên process hợp lệ.

    \item \textbf{Reflective DLL injection:} Tiêm DLL trực tiếp vào
        bộ nhớ mà không cần ghi xuống đĩa, tránh bị antivirus phát
        hiện.

    \item \textbf{Syscall direct invocation:} Mã độc bỏ qua các API
        cấp cao và trực tiếp gọi syscall của Windows, vượt qua hook
        của EDR (Endpoint Detection and Response).

    \item \textbf{Return-oriented programming (ROP):} Dùng các đoạn
        mã ngắn (gadgets) có sẵn trong thư viện để tạo luồng thực thi,
        giúp né DEP và gây khó khăn cho phân tích dòng điều khiển.
\end{itemize}

Việc hiểu rõ đặc điểm và điểm yếu của kiến trúc x86/amd64 và định
dạng PE trên Windows là điều kiện tiên quyết để xây dựng các hệ thống
phát hiện mã độc hiệu quả. Đây cũng là nền tảng để áp dụng các kỹ
thuật như disassembly, symbolic execution, feature extraction và học
máy trong phân tích mã độc hiện đại.















\subsection{Tổng quan về các phương pháp phân tích mã độc PE}

Trước khi đưa vào mô hình học máy để phát hiện và phân loại (là lành
tính hay có chứa mã độc, mã độc thuộc loại nào\dots), một tập tin PE
cần được phân tích và trích xuất đặc trưng. Có ba phương pháp phân
tích tập tin PE cơ bản: phân tích tĩnh, phân tích động và phương pháp
lai. \cite{singh2021survey, damodaran2017comparison}

\textbf{Phân tích tĩnh} (static analysis) trích xuất đặc trưng mà
không cần thực thi mẫu mã độc. Nhiều loại đặc trưng tĩnh khác nhau
có thể được thu thập, bao gồm: giá trị băm (hash), chuỗi N-gram, mã
vận hành (opcode), chuỗi văn bản (string), và thông tin tiêu đề PE
(PE header). Dựa trên những đặc trưng này, các phần mềm phát hiện mã
độc như chương trình chống virus, hệ thống phát hiện xâm nhập (IDS),
v.v. được xây dựng. Để kiểm tra mẫu mã độc ở cấp độ mã máy, các tập
tin thực thi được dịch ngược thành mã hợp ngữ (assembly) thông qua
các công cụ như OllyDbg, IDA Pro, Capstone và WinDbg. Mã hợp ngữ sau
đó được kiểm tra nhằm xác định luồng thực thi, cấu trúc hoặc mẫu hành
vi độc hại, từ đó phục vụ phát hiện các biến thể mới của mã độc hiện
có. Tuy nhiên, phân tích mã hợp ngữ là quá trình tốn thời gian. Hơn
nữa, các kỹ thuật làm rối mã (code obfuscation) - như mã hóa đoạn
mã, hoán đổi thứ tự lệnh, hay chèn mã không thực thi (dead code) -
còn khiến việc phân tích trở nên phức tạp hơn.

\textbf{Phân tích động} (dynamic analysis) thực thi tập tin PE và ghi
lại các hoạt động của chương trình trong thời gian chạy, bao gồm: thao
tác với hệ thống tập tin, thay đổi khóa registry, tạo tiến trình khác,
các hoạt động mạng, lời gọi API hệ điều hành, gửi dữ liệu đến máy chủ
điều khiển (Command and Control - C2)\dots Dựa trên các hoạt động
này, tập tin được phân loại là lành tính (benign) hoặc mã độc
(malicious). Không giống phân tích tĩnh, phân tích động dựa trên sự
tương tác thực tế giữa chương trình và hệ thống. Các tập tin PE chứa
mã độc được thực thi trong môi trường ảo hóa được kiểm soát (sử dụng
các công cụ như VirtualBox hoặc VMware) để tránh gây hại cho hệ thống
thật.

\textbf{Phương pháp lai} (hybrid) kết hợp các
yếu tố của cả phân tích tĩnh và phân tích động. Chẳng hạn, công cụ
``HDM Analyser'' \cite{singh2021survey, damodaran2017comparison} sử dụng cả hai phương pháp trong giai đoạn huấn luyện
(training phase), nhưng chỉ áp dụng phân tích tĩnh trong giai đoạn
kiểm tra (testing phase) - nhằm tận dụng độ chính xác cao của phân
tích động khi huấn luyện, đồng thời duy trì hiệu năng của phân tích
tĩnh khi triển khai. So sánh cho thấy HDM Analyser đạt độ chính xác
tổng thể và độ phức tạp thời gian tốt hơn so với các phương pháp chỉ
dùng một trong hai loại phân tích.

Khó khăn cốt lõi của phương pháp lai là sự bất đối xứng về đặc trưng
khả dụng giữa hai giai đoạn: trong huấn luyện, cả đặc trưng tĩnh lẫn
động đều có thể được thu thập đầy đủ trong môi trường kiểm soát, nhưng trong
triển khai thực tế, đôi khi đặc trưng động không thể được thu thập đầy đủ vì không thể thực
thi tùy tiện một file PE đáng ngờ trên hệ thống thật. Do đó, mô hình
được huấn luyện với đặc trưng phong phú hơn so với những gì có thể
cung cấp tại thời điểm triển khai (inference) - dẫn đến suy giảm hiệu năng
trong vận hành thực tế. Để giải quyết vấn đề trên, có một số phương pháp như sau:

\begin{itemize}
    \item \textbf{Feature distillation:} Huấn luyện mô hình nhận cả
        đặc trưng tĩnh và động, sau đó phân tích các đặc trưng tĩnh
        quan trọng nhất. Phương pháp này có thể gây mất thông tin nếu
        đặc trưng động mới là yếu tố quyết định (có feature importance
        lớn hơn).

    \item \textbf{Pseudo-dynamic features:} Tạo đặc trưng động giả
        bằng cách huấn luyện một mô hình phụ để dự đoán đặc trưng động
        từ đặc trưng tĩnh. Chất lượng dự đoán của kiến trúc tổng thể
        do đó phụ thuộc khá lớn vào độ chính xác của mô hình phụ.

    \item \textbf{Knowledge distillation:} Huấn luyện một mô hình giáo
        viên (teacher) $T$ nhận đầu vào là các đặc trưng tĩnh và động,
        cho đầu ra là logits. Sử dụng $T$ để tạo lại dữ liệu huấn luyện
        chỉ gồm đặc trưng tĩnh và các logits tương ứng, rồi dùng tập
        dữ liệu mới này để huấn luyện mô hình học sinh (student) $S$
        (nhận đầu vào là đặc trưng tĩnh, đầu ra là logits). Mục tiêu
        là $S$ bắt chước đầu ra của $T$ mà không cần đến đặc trưng
        động. Việc dùng logits thay vì nhãn cứng (hard labels) giúp $S$
        phân biệt được các ranh giới ``mềm'' giữa các lớp đầu ra, từ
        đó phân loại tốt hơn.
\end{itemize}

Ba phương pháp trên thể hiện các chiến lược khác nhau nhằm giải quyết
bài toán mất cân bằng giữa khả năng phân biệt cao của đặc trưng động
và yêu cầu hiệu năng, an toàn trong triển khai thực tế - vốn là ưu
điểm của đặc trưng tĩnh. Feature distillation phù hợp với các hệ thống
yêu cầu mô hình đơn giản, dễ diễn giải. Pseudo-dynamic features là lựa
chọn tốt khi cần duy trì cấu trúc phân lớp ban đầu (chẳng hạn giữ
nguyên chiều của vector đầu vào). Knowledge distillation là hướng tiếp
cận tiên tiến hơn, giúp truyền tải hiệu quả phân biệt từ đặc trưng
động sang mô hình gọn nhẹ, phù hợp với triển khai thực tế quy mô lớn.

Việc lựa chọn phương pháp không chỉ là bài toán kỹ thuật, mà còn phản
ánh sự cân bằng giữa độ chính xác, hiệu năng và khả năng mở rộng -
những yếu tố cốt lõi trong thiết kế hệ thống phát hiện mã độc hiện đại.
Trong phạm vi khóa luận, để tập trung vào chất lượng mô hình và khả
năng triển khai thực tế, \textit{phương pháp phân tích tĩnh
kết hợp với trích xuất năng lực (capability extraction)} được
lựa chọn sử dụng. Trong đó, \textit{capability extraction} là kỹ thuật
xác định các hành vi và khả năng tiềm ẩn của file PE thông qua phân
tích tĩnh mã nhị phân, mà không cần thực thi file. Tuy bản chất vẫn
là phân tích tĩnh, capability extraction cho phép thu được thông tin
mang tính ngữ nghĩa hành vi - vốn là ưu điểm đặc trưng của phân
tích động - từ đó hình thành một hướng tiếp cận dung hòa giữa hai phương pháp:
vừa duy trì được hiệu năng và độ an toàn của phân tích tĩnh, vừa bổ
sung chiều sâu ngữ nghĩa mà các kỹ thuật phân tích tĩnh truyền thống
thường không xem xét đầy đủ.










\section{Tổng quan về ứng dụng học máy trong phát hiện mã độc}

Phát hiện mã độc bằng học máy là quá trình sử dụng các thuật toán học
máy để xây dựng mô hình phân loại, cho phép tự động phân biệt giữa
phần mềm hợp lệ và mã độc dựa trên các đặc trưng được trích xuất từ
dữ liệu. Khác với phương pháp dựa trên chữ ký, học máy không phụ
thuộc hoàn toàn vào các mẫu mã độc đã biết, mà có khả năng khái quát
hóa từ dữ liệu huấn luyện để phát hiện các biến thể mới hoặc các mẫu
mã độc chưa từng xuất hiện trước đó.

Trong bối cảnh mã độc PE trên hệ điều hành Windows ngày càng được
thiết kế tinh vi, thường xuyên sử dụng các kỹ thuật đóng gói, làm rối
mã và né tránh phân tích, học máy cho phép khai thác mối quan hệ tiềm
ẩn giữa các đặc trưng mà con người khó có thể xác định bằng các quy
tắc thủ công. Các đặc trưng này có thể xuất phát từ cấu trúc file PE,
hành vi gọi API, hoặc các biểu diễn trừu tượng hơn được học tự động
bởi các mô hình học sâu.

Từ góc độ triển khai, một hệ thống phát hiện mã độc dựa trên học máy
thường tuân theo một quy trình tổng quát, bao gồm: thu thập dữ liệu,
phân tích và trích xuất đặc trưng, huấn luyện mô hình, và đánh giá
hiệu quả. Việc áp dụng quy trình này giúp đảm bảo tính hệ thống, khả
năng tái lập và tính khách quan trong đánh giá kết quả.

Liên quan đến dữ liệu đầu vào, học máy có thể được áp dụng cho nhiều
hướng phân tích mã độc khác nhau, bao gồm phân tích tĩnh, phân tích
động hoặc kết hợp cả hai. Bên cạnh đó, sự phát triển của học sâu đã mở ra khả năng tự động học
biểu diễn đặc trưng từ dữ liệu thô, chẳng hạn như byte sequence hoặc
opcode, thay vì phụ thuộc hoàn toàn vào đặc trưng được thiết kế thủ
công. Các mô hình học sâu tận dụng ưu điểm này để nâng cao khả năng
phát hiện các mẫu mã độc phức tạp, trong khi các mô hình học máy
truyền thống như LightGBM lại thể hiện ưu thế về hiệu năng và tính ổn
định khi triển khai.

Như vậy, phát hiện mã độc bằng học máy không chỉ là sự thay thế đơn
thuần cho các phương pháp truyền thống, mà là một hướng tiếp cận mang
tính tổng hợp, kết hợp kiến thức về hệ điều hành, định dạng PE, kỹ
thuật phân tích mã độc và các thuật toán học máy hiện đại. Các phần
tiếp theo sẽ lần lượt trình bày các hướng tiếp cận phổ biến, kỹ thuật
trích xuất đặc trưng, và mô hình học máy được lựa chọn trong khuôn
khổ khóa luận.

\subsection{Các hướng tiếp cận phổ biến trong phát hiện mã độc bằng học máy}

Trong lĩnh vực phát hiện mã độc, đặc biệt là đối với mã độc định dạng
PE trên hệ điều hành Windows, các phương pháp dựa trên học máy hiện
nay có thể được phân loại theo nhiều hướng tiếp cận khác nhau, tùy
thuộc vào loại dữ liệu đầu vào, cách trích xuất đặc trưng và mô hình
học máy được sử dụng.

\textbf{Hướng tiếp cận dựa trên học máy truyền thống.}
Học máy truyền thống là hướng tiếp cận sớm và vẫn được sử dụng rộng
rãi trong phát hiện mã độc. Các mô hình tiêu biểu trong nhóm này bao
gồm Logistic Regression, Support Vector Machine (SVM), Random Forest,
XGBoost và LightGBM. Đặc điểm chung của các mô hình này là yêu cầu
đặc trưng đầu vào được thiết kế thủ công, dựa trên kiến thức chuyên
gia về cấu trúc file PE và hành vi của mã độc. Hiệu quả của mô hình
do đó phụ thuộc lớn vào chất lượng và tính phân biệt của tập đặc
trưng. Tuy nhiên, các mô hình học máy truyền thống thường có ưu điểm
về thời gian huấn luyện và suy luận nhanh, khả năng diễn giải kết quả
tương đối tốt, và dễ dàng triển khai trong các hệ thống thực tế.

\textbf{Hướng tiếp cận dựa trên học sâu.}
Học sâu (Deep Learning) tận dụng các mạng nơ-ron nhiều tầng để tự
động học biểu diễn đặc trưng từ dữ liệu thô. Trong bài toán phát hiện
mã độc PE, học sâu có thể được áp dụng trên nhiều dạng dữ liệu khác
nhau như byte sequence, opcode, hoặc các biểu diễn trung gian của file
PE. Ưu điểm chính là khả năng phát hiện các mẫu phức tạp và mối quan
hệ phi tuyến trong dữ liệu, đặc biệt hiệu quả khi đối mặt với các kỹ
thuật làm rối và biến thể mã độc. Tuy nhiên, học sâu thường đòi hỏi
lượng dữ liệu huấn luyện lớn, tài nguyên tính toán cao và thời gian
huấn luyện dài hơn so với học máy truyền thống.

\textbf{Nhận xét chung.}
Mỗi hướng tiếp cận trong phát hiện mã độc bằng học máy đều có những
ưu điểm và hạn chế riêng. Việc lựa chọn hướng tiếp cận phù hợp cần
cân nhắc giữa độ chính xác, hiệu năng và khả năng triển khai trong
thực tế. Trong đề tài này, hướng tiếp cận chính là học máy truyền
thống kết hợp phân tích tĩnh, sử dụng mô hình LightGBM làm nền tảng.
Trọng tâm thực nghiệm không nằm ở việc so sánh các loại mô hình, mà
ở việc đánh giá hiệu quả của các tổ hợp đặc trưng khác nhau - cụ
thể là đặc trưng EMBER2024 thuần túy, kết hợp với đặc trưng năng lực
(capability) trích xuất bằng công cụ Capa ở dạng one-hot và dạng đã
giảm chiều bằng TruncatedSVD - nhằm xác định tổ hợp đặc trưng tối
ưu cho bài toán phát hiện mã độc PE.






\subsection{Trích xuất đặc trưng từ file PE}

Trong bài toán phát hiện mã độc dựa trên học máy, quá trình trích xuất
đặc trưng đóng vai trò then chốt, quyết định trực tiếp đến hiệu quả
của mô hình. Đối với các file thực thi định dạng PE, các đặc trưng có
thể được trích xuất từ nhiều khía cạnh khác nhau, phản ánh cấu trúc,
nội dung và một phần hành vi tiềm ẩn của chương trình. Phần này trình
bày các nhóm đặc trưng chính được sử dụng trong đề tài, đặc biệt là
trong các feature extractor (bộ trích xuất đặc trưng) nổi tiếng như EMBERv2 và EMBER2024
\cite{ember2024benchmark}.

\subsubsection{Đặc trưng tổng quát của file PE}

Nhóm đặc trưng này phản ánh các thông tin tổng quát của file PE, bao
gồm kích thước file, số lượng section, thời gian biên dịch (timestamp) và các cờ
(flag) cấu hình trong header. Các đặc trưng này giúp mô hình học máy
nắm bắt được những khác biệt cơ bản giữa file thực thi hợp lệ và mã
độc, bởi mã độc thường có cấu trúc bất thường hoặc được tạo ra thông
qua các công cụ tự động như packer, obfuscator, cryptor\dots

\subsubsection{Đặc trưng cấu trúc PE và header}

Các đặc trưng cấu trúc được trích xuất từ các thành phần quan trọng
trong file PE như DOS header, PE header và optional header. Nhóm đặc
trưng này bao gồm các trường thông tin phản ánh cách thức file được
nạp và thực thi trên hệ điều hành Windows. Những sai lệch hoặc giá
trị bất thường trong các trường header thường là dấu hiệu của mã độc
hoặc các file đã bị chỉnh sửa.

\subsubsection{Đặc trưng liên quan đến section}

Mỗi file PE được chia thành nhiều section với các mục đích khác nhau.
Các đặc trưng liên quan đến section bao gồm số lượng section, kích
thước của từng section, quyền truy cập (read, write, execute) và phân
bố entropy giữa các section. Đặc biệt, entropy cao ở một số section
có thể là dấu hiệu của kỹ thuật đóng gói hoặc mã hóa, vốn thường
được sử dụng trong mã độc để né tránh phân tích.

Trong lý thuyết thông tin, entropy được sử dụng như một thước đo mức
độ hỗn loạn và tính ngẫu nhiên của dữ liệu. Đối với các chương trình
lành tính thông thường, dữ liệu trong các section của file thực thi
thường có cấu trúc tương đối rõ ràng; khi tiến hành dịch ngược, có
thể quan sát được luồng điều khiển và logic chương trình ở mức cơ bản.
Ngược lại, nhiều loại mã độc sử dụng các kỹ thuật đóng gói hoặc mã
hóa, làm cho nội dung các section trở nên khó phân tích và có mức
entropy cao hơn so với các file thực thi thông thường.

Tuy nhiên, cần lưu ý rằng entropy cao không phải là dấu hiệu đặc
trưng tuyệt đối của mã độc. Một số phần mềm thương mại cũng áp dụng
các kỹ thuật bảo vệ mã nguồn tương tự, dẫn đến mức entropy cao ở một
số section. Do đó, việc đánh giá mã độc chỉ dựa trên entropy là không
đủ và cần kết hợp với các nhóm đặc trưng khác để đưa ra quyết định
phân loại chính xác hơn.

\subsubsection{Đặc trưng về bảng import và export}

Nhóm đặc trưng này tập trung vào các hàm và thư viện được file PE sử
dụng. Việc phân tích các thư viện và hàm import cho phép suy đoán hành
vi tiềm ẩn của chương trình, chẳng hạn như thao tác với hệ thống file,
registry, mạng hoặc tiến trình. Các file mã độc thường sử dụng các hàm
API đặc trưng phục vụ cho việc ẩn mình, duy trì tồn tại hoặc thu thập
thông tin.

\subsubsection{Đặc trưng thống kê từ nội dung nhị phân}

Nhóm đặc trưng thống kê được trích xuất trực tiếp từ nội dung nhị phân
của file, bao gồm phân bố byte, tần suất xuất hiện của các giá trị
byte và các thống kê liên quan. Những đặc trưng này phản ánh cấu trúc
tổng thể của file ở mức thấp và có khả năng phân biệt giữa file hợp
lệ và mã độc ngay cả khi mã đã bị làm rối.

\subsubsection{Đặc trưng trích xuất từ chuỗi văn bản}

Trong file PE, nhiều chuỗi văn bản có thể được trích xuất, chẳng hạn
như tên file, đường dẫn, URL, tên hàm hoặc thông báo lỗi. Việc phân
tích các chuỗi này giúp phát hiện các dấu hiệu liên quan đến hoạt
động độc hại, như kết nối mạng, tải thêm mã hoặc tương tác với hệ
thống. Các chuỗi văn bản thường được mã độc sử dụng để phục vụ cho
quá trình thực thi hoặc giao tiếp với máy chủ điều khiển.

\subsubsection{Đặc trưng dạng biểu diễn rút gọn}

Để phù hợp với các mô hình học máy truyền thống, nhiều đặc trưng có
tính chất danh mục hoặc chuỗi được chuyển đổi sang dạng biểu diễn số
thông qua các kỹ thuật rút gọn và mã hóa. Các kỹ thuật này giúp giảm
chiều dữ liệu, đồng thời vẫn giữ lại các thông tin quan trọng phục vụ
cho quá trình học của mô hình.

\subsubsection{Đặc trưng dựa trên chữ ký số}

Trong hệ điều hành Windows, các file thực thi PE có thể được ký số
nhằm xác thực nguồn gốc và đảm bảo tính toàn vẹn của file. Các đặc
trưng liên quan đến chữ ký số phản ánh trạng thái xác thực của file,
chẳng hạn như việc file có được ký số hay không, chữ ký có hợp lệ hay
không, cũng như một số thuộc tính cơ bản của chứng thư số đi kèm.

Trong thực tế, phần lớn các phần mềm hợp lệ thường được ký số để tăng
độ tin cậy, trong khi nhiều mẫu mã độc không được ký số hoặc sử dụng
chữ ký không hợp lệ, tự ký hoặc đã bị chỉnh sửa. Tuy nhiên, sự tồn
tại của chữ ký số hợp lệ không đảm bảo tuyệt đối file là an toàn,
bởi mã độc vẫn có thể được phát tán thông qua các chứng thư bị lạm
dụng hoặc bị đánh cắp. Vì vậy, nhóm đặc trưng này được sử dụng như
một yếu tố bổ trợ, kết hợp với các nhóm đặc trưng khác trong quá
trình phân loại.

\subsubsection{Đặc trưng tổng hợp các lỗi đọc file PE}

Quá trình phân tích tĩnh file PE đòi hỏi việc đọc và diễn giải chính
xác cấu trúc nội bộ của file theo định dạng chuẩn. Trong thực tế,
nhiều file mã độc được tạo ra với cấu trúc bất thường, bị làm hỏng có
chủ đích hoặc cố tình vi phạm chuẩn định dạng nhằm gây khó khăn cho
các công cụ phân tích. Do đó, các lỗi phát sinh khi đọc hoặc phân
tích cấu trúc file PE - như header không hợp lệ, bảng section bị
thiếu hoặc trùng lặp, offset không nhất quán, hoặc các trường dữ liệu
chứa giá trị bất thường - có thể được ghi nhận và tổng hợp thành
một nhóm đặc trưng riêng.

Về mặt thống kê, một file thực thi phát sinh càng nhiều lỗi trong quá
trình phân tích thì mức độ bất thường của file càng cao, thường phản
ánh việc file đã bị chỉnh sửa, làm rối hoặc được tạo ra bởi các công
cụ tự động - những đặc điểm phổ biến của mã độc. Tuy nhiên, tương
tự như các nhóm đặc trưng khác, lỗi phân tích không thể được sử dụng
độc lập để kết luận một file là mã độc, bởi một số file hợp lệ cũng
có thể bị hỏng hoặc không tuân thủ hoàn toàn chuẩn định dạng. Vì vậy,
nhóm đặc trưng này được kết hợp với các đặc trưng cấu trúc, thống kê
và nội dung khác nhằm nâng cao độ chính xác của mô hình.

Ngoài các nhóm đặc trưng nêu trên, một số hướng mở rộng có thể được
xem xét trong các nghiên cứu tiếp theo, bao gồm đặc trưng dựa trên
opcode hoặc lệnh máy (assembly), đặc trưng dựa trên đồ thị luồng điều
khiển, và đặc trưng kết hợp giữa phân tích tĩnh và động. Tuy nhiên,
trong phạm vi đề tài, các nhóm đặc trưng tĩnh từ file PE được lựa
chọn nhằm đảm bảo tính an toàn, khả năng mở rộng và phù hợp với yêu
cầu thực nghiệm.

\subsection{Một số mô hình học máy tiêu biểu trong phát hiện mã độc PE}

Trong những năm gần đây, nhiều mô hình học máy đã được nghiên cứu và
áp dụng vào bài toán phát hiện mã độc, đặc biệt là đối với các file
thực thi định dạng PE trên hệ điều hành Windows. Tùy thuộc vào đặc
điểm dữ liệu, phương pháp trích xuất đặc trưng và yêu cầu triển khai,
các mô hình học máy khác nhau thể hiện những ưu điểm và hạn chế riêng.

\textbf{Hồi quy Logistic} là một trong những mô hình học máy đơn giản
và phổ biến nhất cho bài toán phân loại nhị phân. Trong phát hiện mã
độc, mô hình này được sử dụng để ước lượng xác suất một file thực thi
là mã độc dựa trên các đặc trưng đầu vào. Ưu điểm của hồi quy
Logistic là dễ cài đặt, tốc độ huấn luyện nhanh và khả năng diễn giải
cao. Tuy nhiên, do giả định mối quan hệ tuyến tính giữa đặc trưng và
nhãn, mô hình này thường gặp hạn chế khi xử lý dữ liệu phức tạp và
có tính phi tuyến cao như trong các tập dữ liệu mã độc thực tế.

\textbf{Support Vector Machine (SVM)} là mô hình học máy mạnh mẽ trong
việc phân loại dữ liệu có số chiều lớn. SVM hoạt động bằng cách tìm
siêu phẳng tối ưu nhằm phân tách các lớp dữ liệu với khoảng cách lớn
nhất. Trong bài toán phát hiện mã độc, SVM cho kết quả tương đối tốt
và tỏ ra vượt trội hơn hồi quy logistic khi sử dụng các kernel phi
tuyến. Tuy nhiên, chi phí tính toán của SVM có thể tăng cao khi kích
thước tập dữ liệu lớn, khiến việc triển khai trong môi trường thực tế
gặp nhiều hạn chế.

\textbf{Decision Tree và Random Forest.} Cây quyết định (Decision Tree)
là mô hình học máy dựa trên cấu trúc cây, trong đó mỗi nút đại diện
cho một điều kiện phân tách dữ liệu. Mô hình này có ưu điểm là trực
quan, dễ hiểu và có khả năng xử lý dữ liệu phi tuyến. Rừng ngẫu nhiên
(Random Forest) là mô hình mở rộng của cây quyết định, kết hợp nhiều
cây độc lập để nâng cao độ chính xác và giảm hiện tượng quá khớp.
Trong phát hiện mã độc, rừng ngẫu nhiên thường được sử dụng nhờ khả
năng xử lý tốt dữ liệu nhiều chiều và tính ổn định cao.

\textbf{Các mô hình tăng cường dựa trên cây} như Gradient Boosting,
XGBoost và LightGBM đã cho thấy hiệu quả vượt trội trong nhiều bài
toán phân loại, bao gồm phát hiện mã độc. Các mô hình này xây dựng
tập hợp các cây yếu theo cách tuần tự, trong đó mỗi cây mới tập trung
khắc phục các sai sót của các cây trước đó. Ưu điểm của nhóm mô hình
này là khả năng học các mối quan hệ phi tuyến phức tạp, xử lý hiệu
quả dữ liệu có số chiều lớn và khả năng kiểm soát hiện tượng quá khớp
thông qua các tham số điều chỉnh. Trong đề tài này, LightGBM được lựa
chọn để triển khai và đánh giá chi tiết hơn.

Các mô hình học máy truyền thống, đặc biệt là các mô hình dựa trên
cây và tăng cường, vẫn đóng vai trò quan trọng trong phát hiện mã độc
nhờ sự cân bằng giữa độ chính xác, hiệu năng và khả năng triển khai.
Việc lựa chọn mô hình phù hợp phụ thuộc vào đặc điểm tập dữ liệu,
nhóm đặc trưng được sử dụng và yêu cầu cụ thể của bài toán.

\subsection{Mô hình LightGBM}

Light Gradient Boosting Machine (LightGBM) là một mô hình học máy
thuộc nhóm các thuật toán tăng cường dựa trên cây quyết định, được
thiết kế nhằm nâng cao hiệu quả huấn luyện và khả năng mở rộng đối
với các tập dữ liệu có kích thước lớn và số lượng đặc trưng cao. Trong
bối cảnh bài toán phát hiện mã độc PE, LightGBM cho thấy nhiều ưu
điểm phù hợp với yêu cầu thực nghiệm và triển khai.

LightGBM xây dựng mô hình bằng cách kết hợp nhiều cây quyết định yếu
theo cơ chế học tăng cường, trong đó mỗi cây mới được huấn luyện nhằm
giảm sai số còn tồn tại của các cây trước đó. Khác với các phương pháp
tăng cường truyền thống, LightGBM áp dụng chiến lược tăng trưởng cây
theo chiều sâu (leaf-wise), giúp mô hình hội tụ nhanh hơn và khai thác
hiệu quả các mối quan hệ phi tuyến trong dữ liệu.

Một ưu điểm quan trọng của LightGBM là khả năng xử lý hiệu quả dữ
liệu có số chiều lớn, vốn là đặc trưng phổ biến trong bài toán phát
hiện mã độc khi sử dụng nhiều nhóm đặc trưng trích xuất từ file PE.
Bên cạnh đó, LightGBM hỗ trợ cơ chế rời rạc hóa đặc trưng và tối ưu
bộ nhớ, giúp giảm chi phí tính toán trong quá trình huấn luyện và suy
luận.

Trong phạm vi đề tài, LightGBM được lựa chọn làm mô hình đại diện cho
hướng tiếp cận học máy truyền thống nhờ các đặc điểm sau:

\begin{itemize}
    \item Khả năng đạt độ chính xác cao với dữ liệu đặc trưng tĩnh
        từ file PE.
    \item Thời gian huấn luyện và suy luận nhanh, phù hợp với kịch
        bản xử lý số lượng lớn file.
    \item Tính ổn định và khả năng kiểm soát hiện tượng quá khớp
        thông qua các tham số điều chỉnh.
\end{itemize}

Ngoài ra, LightGBM cho phép đánh giá mức độ đóng góp của từng đặc
trưng thông qua các chỉ số feature importance hoặc kỹ thuật xAI
như SHAP, giúp hỗ trợ phân tích
và giải thích kết quả thực nghiệm. Khả năng này đặc biệt hữu ích
trong việc đánh giá vai trò của từng nhóm đặc trưng, cung cấp cơ sở
lý giải cho kết quả thực nghiệm thu được.









\section{Các công trình nghiên cứu liên quan}

Phần này điểm qua các công trình nghiên cứu và tập dữ liệu tiêu biểu
liên quan đến bốn khía cạnh chính của khóa luận: tập dữ liệu phục vụ
huấn luyện mô hình, ứng dụng học máy trong phân tích tĩnh file PE,
kỹ thuật trích xuất năng lực mã độc, và phương pháp lý giải quyết
định của mô hình. Trên cơ sở đó, khóa luận xác định điểm kế thừa và
đóng góp riêng so với các nghiên cứu hiện có.

\subsection{Các tập dữ liệu file PE phục vụ huấn luyện mô hình học máy phát hiện mã độc}

Chất lượng và quy mô tập dữ liệu có ảnh hưởng trực tiếp đến hiệu quả
huấn luyện và khả năng tổng quát hóa của mô hình học máy trong bài
toán phát hiện mã độc. Một số tập dữ liệu tiêu biểu đã được công bố
và sử dụng rộng rãi trong cộng đồng nghiên cứu.

\textbf{EMBER (Endgame Malware BEnchmark for Research).}
EMBER \cite{ember2018} là một trong những tập dữ liệu benchmark được
sử dụng phổ biến nhất cho bài toán phát hiện mã độc PE. Phiên bản đầu
tiên được Anderson và Roth công bố năm 2018, bao gồm khoảng 1,1 triệu
mẫu file PE (lành tính, mã độc và chưa gán nhãn), với các đặc trưng
tĩnh được trích xuất sẵn theo định dạng vector có cấu trúc. Tập dữ
liệu này cung cấp một pipeline trích xuất đặc trưng chuẩn hóa, bao
gồm các nhóm đặc trưng từ PE header, section, import table, chuỗi văn
bản và nội dung nhị phân, giúp tạo điều kiện thuận lợi cho việc so
sánh các mô hình học máy trên cùng một nền tảng dữ liệu.

Phiên bản EMBER2024 \cite{ember2024benchmark} là bản cập nhật mới nhất, được
xây dựng nhằm phản ánh bối cảnh mã độc hiện đại hơn so với phiên bản
2018. EMBER2024 bổ sung số lượng mẫu lớn hơn, cập nhật phân bố thời
gian thu thập, và mở rộng pipeline trích xuất đặc trưng. Trong đề tài
này, EMBER2024 được lựa chọn làm tập dữ liệu chính cho tất cả các
thực nghiệm.

\textbf{VirusShare và VirusTotal.}
VirusShare \cite{virusshare} là một kho lưu trữ mẫu mã độc lớn được
duy trì bởi cộng đồng nghiên cứu bảo mật, cung cấp hàng triệu mẫu
file thực thi đã được xác nhận là mã độc. VirusTotal \cite{virustotal}
là dịch vụ phân tích file trực tuyến tổng hợp kết quả từ hàng chục
engine antivirus, thường được dùng để gán nhãn hoặc xác minh nhãn cho
các mẫu trong quá trình xây dựng tập dữ liệu. Nhiều nghiên cứu sử
dụng kết hợp hai nguồn này để thu thập mẫu mã độc đa dạng và đảm bảo
độ tin cậy của nhãn phân loại.

\textbf{Các tập dữ liệu khác.}
Ngoài EMBER, một số tập dữ liệu khác cũng được sử dụng trong các
nghiên cứu liên quan, chẳng hạn như SOREL-20M \cite{sorel20m} -
tập dữ liệu quy mô lớn với 20 triệu mẫu PE do Sophos và ReversingLabs
công bố - cũng như các tập dữ liệu tự thu thập từ honeypot hoặc các nguồn
chia sẻ trong cộng đồng bảo mật. Tuy nhiên, EMBER vẫn là lựa chọn
phổ biến nhất nhờ tính chuẩn hóa, khả năng tái lập và sự hỗ trợ,
công nhận rộng rãi của cộng đồng nghiên cứu.

\subsection{Các nghiên cứu về ứng dụng học máy trong phân tích tĩnh file thực thi nhị phân}

Ứng dụng học máy trong phân tích tĩnh file PE là một hướng nghiên cứu
được quan tâm rộng rãi, với nhiều công trình tiêu biểu đã đạt kết quả
đáng chú ý.
 
\textbf{Các mô hình học máy truyền thống trên đặc trưng tĩnh.}
Anderson và Roth \cite{ember2018} đã đề xuất mô hình LightGBM huấn
luyện trên tập đặc trưng EMBER, đạt tỷ lệ phát hiện mã độc cao với
tốc độ suy luận phù hợp cho triển khai thực tế. Đây là một trong
những baseline quan trọng nhất trong lĩnh vực, cho thấy các mô hình
tăng cường dựa trên cây quyết định có thể đạt hiệu quả cạnh tranh
với các phương pháp phức tạp hơn khi được kết hợp với bộ đặc trưng
tĩnh chất lượng cao.

\textbf{Kỹ thuật giảm chiều trong xử lý đặc trưng.}
Khi tập đặc trưng có số chiều lớn - đặc biệt khi kết hợp nhiều
nhóm đặc trưng khác nhau - các kỹ thuật giảm chiều như PCA
(Principal Component Analysis) \cite{pca} và TruncatedSVD
\cite{truncatedsvd} thường được áp dụng để giảm chi phí tính toán và
hạn chế hiện tượng quá khớp. TruncatedSVD đặc biệt phù hợp với các
ma trận đặc trưng thưa (sparse), vốn xuất hiện phổ biến khi mã hóa
đặc trưng dạng one-hot hoặc bag-of-words từ dữ liệu mã độc.

\textbf{Phương pháp Ensemble.}
Các phương pháp kết hợp mô hình (ensemble) như voting, stacking và
blending \cite{ensemble_survey} đã được áp dụng trong phát hiện mã
độc nhằm kết hợp điểm mạnh của nhiều mô hình thành phần khác loại
(heterogeneous ensemble), từ đó nâng cao độ chính xác và tính ổn định
tổng thể. Tuy nhiên, cần phân biệt rõ hai cấp độ ensemble: ensemble
giữa các mô hình khác loại (ví dụ kết hợp LightGBM với SVM hay mạng
neuron) và ensemble nội tại bên trong một mô hình. LightGBM thuộc
nhóm thứ hai - bản thân mô hình đã là một tổ hợp của hàng trăm cây
quyết định được xây dựng tuần tự theo cơ chế gradient boosting, trong
đó mỗi cây mới tập trung khắc phục sai số của các cây trước. Cơ chế
này mang lại khả năng học các mối quan hệ phi tuyến phức tạp và tính
ổn định cao, mà không cần thêm một lớp ensemble bên ngoài. Các nghiên
cứu thực nghiệm \cite{ember2018, sorel20m} cho thấy LightGBM đơn lẻ
đã đạt hiệu năng cạnh tranh với các heterogeneous ensemble phức tạp
hơn trong bài toán phát hiện mã độc PE, trong khi vẫn duy trì ưu thế
về tốc độ suy luận và khả năng triển khai thực tế - những yếu tố
đặc biệt quan trọng khi triển khai mô hình ngay trong
dịch vụ chạy trên máy endpoint trong khuôn khổ khóa luận.

\subsection{Các nghiên cứu về trích xuất năng lực (capability extraction) của file thực thi nhị phân}

Trích xuất năng lực (capability extraction) là kỹ thuật xác định các
hành vi và khả năng tiềm ẩn của một file thực thi thông qua phân tích
tĩnh, mà không cần thực thi file trong môi trường thật. Đây là hướng
tiếp cận có tiềm năng cao trong việc bổ sung thông tin ngữ nghĩa hành
vi vào quá trình phát hiện mã độc.

\textbf{Công cụ Capa.}
Capa \cite{capa} là công cụ mã nguồn mở do MANDIANT (trước đây là
FireEye) phát triển, cho phép tự động nhận diện các năng lực của file
thực thi nhị phân thông qua một tập quy tắc (rule) được xây dựng thủ
công bởi các chuyên gia phân tích mã độc. Mỗi quy tắc trong Capa mô
tả một năng lực cụ thể - chẳng hạn ``mã hóa dữ liệu bằng XOR'',
``tạo tiến trình mới'', hay ``kết nối đến địa chỉ IP từ xa'' - và
được ánh xạ sang các framework phân loại hành vi chuẩn như MITRE
ATT\&CK \cite{attack} và MBC (Malware Behavior Catalog) \cite{mbc}.
Capa hỗ trợ phân tích file PE thông qua kỹ thuật phân tích tĩnh 
disassembly (dịch ngược hợp ngữ), sử dụng các backend như SMDA
hoặc IDA Pro để dịch ngược và đối chiếu với tập quy tắc.
 
Trong đề tài này, Capa được sử dụng để trích xuất vector đặc trưng
năng lực từ các file PE trong tập EMBER2024, sau đó kết hợp với đặc
trưng EMBER2024 gốc để huấn luyện mô hình học máy. Điểm khác biệt
so với cách dùng Capa truyền thống (chủ yếu phục vụ phân tích thủ
công) là đề tài tự động hóa quá trình này ở quy mô lớn và tích hợp
kết quả vào pipeline học máy.

Một số nghiên cứu đã khảo sát việc sử dụng thông tin về năng lực và
hành vi của mã độc như một nguồn đặc trưng bổ sung cho mô hình học
máy \cite{kim2022mapas,li2024malapi}. Các kết quả ban đầu cho thấy đặc trưng năng lực
có tính phân biệt cao, đặc biệt trong việc phân loại họ mã độc và
phát hiện các biến thể mới sử dụng kỹ thuật obfuscation. Tuy nhiên,
việc tích hợp đặc trưng Capa vào pipeline học máy ở quy mô lớn vẫn
còn ít được nghiên cứu một cách có hệ thống, đặc biệt khi kết hợp
với các kỹ thuật giảm chiều như TruncatedSVD - đây là một trong
những đóng góp chính của khóa luận.

\subsection{Các nghiên cứu về lý giải quyết định của mô hình (Explainable AI - xAI)}

Khi các mô hình học máy ngày càng được triển khai trong các hệ thống
bảo mật thực tế, khả năng lý giải quyết định của mô hình (Explainable
AI - xAI) trở thành yêu cầu quan trọng, giúp chuyên gia bảo mật
hiểu được lý do một file bị phân loại là mã độc và đưa ra quyết định
phản ứng phù hợp.
 
\textbf{SHAP (SHapley Additive exPlanations).}
SHAP \cite{shap} là một trong những phương pháp xAI phổ biến và có
nền tảng lý thuyết vững chắc nhất hiện nay, dựa trên khái niệm Shapley
value từ lý thuyết trò chơi hợp tác. SHAP tính toán mức độ đóng góp
của từng đặc trưng vào kết quả dự đoán của mô hình cho từng mẫu cụ
thể, cho phép giải thích quyết định ở cấp độ cục bộ (per-instance)
lẫn toàn cục (global feature importance). Trong bài toán phát hiện mã
độc, SHAP được sử dụng để xác định những đặc trưng nào - chẳng hạn
một năng lực Capa cụ thể hay một trường PE header - đóng vai trò
quyết định trong việc phân loại một file là mã độc.
 
\textbf{LIME (Local Interpretable Model-agnostic Explanations).}
LIME \cite{lime} là phương pháp xAI không phụ thuộc vào loại mô hình
(model-agnostic), hoạt động bằng cách xây dựng một mô hình tuyến tính
đơn giản xấp xỉ hành vi của mô hình phức tạp trong vùng lân cận của
một điểm dữ liệu cụ thể. LIME phù hợp với các trường hợp cần giải
thích nhanh và trực quan cho từng mẫu riêng lẻ, mặc dù độ ổn định
của giải thích có thể thấp hơn so với SHAP trong một số trường hợp.
 
\textbf{Ứng dụng xAI trong phát hiện mã độc.}
Một số nghiên cứu đã kết hợp SHAP với các mô hình phát hiện mã độc
dựa trên học máy \cite{shap_malware}, cho thấy xAI không chỉ hỗ trợ
kiểm tra và gỡ lỗi mô hình mà còn cung cấp thông tin hữu ích cho quá
trình điều tra và ứng phó sự cố. Trong đề tài này, SHAP được tích hợp
vào server phân tích sâu (deep-analysis-server) nhằm cung cấp giải
thích trực quan cho kết quả phân loại, giúp người dùng hiểu được
những đặc trưng nào - bao gồm cả đặc trưng Capa - đã dẫn đến
quyết định phân loại của mô hình.

\textbf{Nhận xét chung.}
So với các công trình liên quan, đề tài này có những điểm kế thừa và
khác biệt như sau. Về kế thừa, đề tài sử dụng tập dữ liệu EMBER2024
và mô hình LightGBM theo hướng đã được thiết lập tốt trong cộng đồng
nghiên cứu, đồng thời áp dụng SHAP - một kỹ thuật xAI đã được kiểm
chứng - để giải thích kết quả. Về đóng góp riêng, đề tài tập trung
vào việc tích hợp có hệ thống đặc trưng năng lực từ Capa vào pipeline
học máy ở quy mô lớn, đánh giá định lượng mức độ đóng góp của đặc
trưng này so với baseline EMBER2024 thuần túy, và triển khai
pipeline vào một dịch vụ bảo vệ endpoint thực tế trên Windows.
 









\section{Công nghệ sử dụng}

Đề tài kết hợp nhiều công nghệ và công cụ thuộc các lĩnh vực khác
nhau, từ học máy và phân tích mã độc đến lập trình hệ thống Windows.
Phần này trình bày tổng quan các công nghệ chính được sử dụng trong
quá trình thực hiện.

\subsection{Ngôn ngữ lập trình và công cụ xây dựng}

\textbf{Python} là ngôn ngữ lập trình chính cho các thành phần liên
quan đến học máy và phân tích dữ liệu. Python được lựa chọn nhờ hệ
sinh thái thư viện phong phú, đặc biệt trong lĩnh vực khoa học dữ
liệu và học máy, cùng với khả năng tích hợp nhanh chóng với các công
cụ phân tích mã độc hiện có.

\textbf{C++} được sử dụng cho module trích xuất đặc trưng EMBER2024
chạy trực tiếp trên máy endpoint. Lựa chọn này xuất phát từ yêu cầu
về hiệu năng - việc trích xuất đặc trưng cần được thực hiện nhanh
chóng trên từng file PE trong thời gian thực - cũng như khả năng
tích hợp tự nhiên với các API hệ thống Windows. Dự án C++ được quản
lý bằng \textbf{CMake}, công cụ build đa nền tảng phổ biến, giúp đảm
bảo tính nhất quán trong quá trình biên dịch và triển khai.

\subsection{Lập trình hệ thống Windows}

Các thành phần chạy trực tiếp trên máy endpoint được xây dựng dựa
trên các công nghệ lập trình hệ thống của Windows.

\textbf{Windows API} \cite{winapi} là tập hợp các hàm và giao diện
lập trình do Microsoft cung cấp, cho phép ứng dụng tương tác trực
tiếp với hệ điều hành Windows ở mức người dùng (user mode). Trong đề
tài, Windows API được sử dụng để xây dựng dịch vụ nền (background
service) và giao diện người dùng (UI) chạy trên endpoint, bao gồm
các chức năng quản lý tiến trình, giao tiếp liên tiến trình (IPC),
tương tác với hệ thống file, cùng một số chức năng khác.

\textbf{Kernel-Mode Driver Framework (KMDF)} \cite{kmdf} là framework
do Microsoft cung cấp để phát triển driver chạy ở chế độ nhân
(kernel mode) trên Windows. KMDF cung cấp một lớp trừu tượng trên
WDM (Windows Driver Model), giúp đơn giản hóa việc lập trình driver
và giảm thiểu rủi ro gây mất ổn định hệ thống. Trong đề tài, KMDF
được sử dụng để xây dựng thành phần driver có khả năng theo dõi sự
kiện thực thi file ở tầng nhân, phục vụ cho cơ chế phát hiện mã độc
theo thời gian thực.

\subsection{Phát triển dịch vụ phân tích}

\textbf{FastAPI} \cite{fastapi} là framework Python hiện đại để xây
dựng REST API, nổi bật với hiệu năng cao nhờ hỗ trợ bất đồng bộ
(async) và khả năng tự động sinh tài liệu API. FastAPI được sử dụng
làm nền tảng cho server phân tích sâu, tiếp nhận yêu cầu phân tích
file từ các thành phần khác trong hệ thống và trả về kết quả phân
loại cùng giải thích từ mô hình.

Giao diện web của server phân tích được xây dựng bằng \textbf{HTML,
CSS và JavaScript} thuần, không phụ thuộc vào framework frontend phức
tạp, nhằm đảm bảo tính đơn giản và dễ triển khai.

\textbf{Docker} \cite{docker} được sử dụng để đóng gói và triển khai
server phân tích sâu dưới dạng container, đảm bảo tính nhất quán về
môi trường thực thi và đơn giản hóa quá trình cài đặt trên các máy
khác nhau.

\subsection{Thư viện học máy và phân tích mã độc}

\textbf{LightGBM} \cite{lgbm} là thư viện gradient boosting hiệu năng
cao do Microsoft phát triển, được sử dụng làm mô hình phân loại chính
trong đề tài. LightGBM hỗ trợ huấn luyện song song, tối ưu bộ nhớ
và cung cấp API tích hợp tốt với hệ sinh thái Python. Ngoài ra, thư
viện còn có phiên bản (binding) ngôn ngữ C++, hỗ trợ việc triển
khai mô hình trực tiếp trên máy endpoint, đảm bảo hiệu năng mà
không cần Python runtime.

\textbf{scikit-learn} \cite{sklearn} là thư viện học máy nền tảng
trong Python, được sử dụng cho các tác vụ tiền xử lý dữ liệu, giảm
chiều đặc trưng (TruncatedSVD) và đánh giá mô hình.

\textbf{SHAP} \cite{shap} là thư viện Python cung cấp các phương pháp
giải thích mô hình học máy dựa trên Shapley value. SHAP được tích hợp
vào server phân tích sâu để tạo ra các giải thích trực quan về lý do
mô hình phân loại một file là mã độc.

\textbf{Capa} \cite{capa} là công cụ mã nguồn mở do Mandiant phát
triển, cho phép tự động nhận diện năng lực (capability) của file thực
thi nhị phân thông qua phân tích tĩnh dựa trên tập quy tắc. Trong đề
tài, Capa được tích hợp vào pipeline trích xuất đặc trưng để bổ sung
thông tin ngữ nghĩa hành vi vào quá trình huấn luyện mô hình, cũng
như được sử dụng trong dịch vụ phân tích chuyên sâu online
(deep-analysis-server).
